%
% hybrid.tex -- Übergang Quadtree -> Newton
%
% (c) 2021 Prof Dr Andreas Müller, OST Ostschweizer Fachhochschule
%
\documentclass[tikz]{standalone}
\usepackage{amsmath}
\usepackage{times}
\usepackage{txfonts}
\usepackage{pgfplots}
\usepackage{csvsimple}
\usetikzlibrary{arrows,intersections,math}
\begin{document}
\def\skala{1.08}
\begin{tikzpicture}[>=latex,thick,scale=\skala,
    bigcell/.style={draw=blue!50!black,fill=blue!10,line width=0.6pt,rounded corners=1pt},
    smallcell/.style={draw=orange!70!black,fill=orange!20,line width=0.5pt,rounded corners=0.5pt},
    isocell/.style={draw=teal!60!black,fill=teal!20,line width=0.6pt,rounded corners=1pt},
    root/.style={circle,fill=red!80!black,inner sep=0pt,minimum size=3.5pt},
    newtonpt/.style={circle,fill=teal!60!black,inner sep=0pt,minimum size=2.5pt},
    label/.style={font=\bfseries},
    sub/.style={gray!50!black}]
%
% =========== LINKE SEITE: Quadtree fast fertig ===========
\begin{scope}[xshift=0cm]
  \node[label] at (2,4.4) {Quadtree-Phase};
  \node[sub] at (2,4.0) {grobe Lokalisation};
  % Rahmen
  \draw[gray!30,thin] (0,0) rectangle (4,3.7);
  % Mehrere kleine Quadrate, jedes mit einer Nullstelle
  \draw[isocell] (0.3,2.7) rectangle (0.9,3.3);
  \node[root] at (0.6,3.0) {};
  \node[sub] at (0.6,2.5) {$N{=}1$};
  %
  \draw[isocell] (2.7,3.0) rectangle (3.3,3.6);
  \node[root] at (3.0,3.3) {};
  \node[sub] at (3.0,2.8) {$N{=}1$};
  %
  \draw[isocell] (1.5,1.7) rectangle (2.1,2.3);
  \node[root] at (1.8,2.0) {};
  \node[sub] at (1.8,1.5) {$N{=}1$};
  %
  \draw[isocell] (3.1,1.2) rectangle (3.7,1.8);
  \node[root] at (3.4,1.5) {};
  \node[sub] at (3.4,1.0) {$N{=}1$};
  %
  \draw[isocell] (0.5,0.5) rectangle (1.1,1.1);
  \node[root] at (0.8,0.8) {};
  \node[sub] at (0.8,0.3) {$N{=}1$};
\end{scope}
%
% Pfeil mit Übergangs-Label
\draw[->,line width=1.0pt,gray!40!black] (4.4,1.9) -- (5.5,1.9);
\node[gray!40!black] at (4.95,2.3) {Übergabe};
\node[gray!50!black] at (4.95,1.5) {Mittelpunkte};
%
% =========== RECHTE SEITE: Newton-Konvergenz ===========
\begin{scope}[xshift=6cm]
  \node[label] at (2,4.4) {Newton-Phase};
  \node[sub] at (2,4.0) {schnelle Feinkonvergenz};
  \draw[gray!30,thin] (0,0) rectangle (4,3.7);
  %
  % Ein vergrössertes Quadrat aus dem Quadtree, das eine Nullstelle enthält
  \draw[isocell] (1.0,1.0) rectangle (3.0,3.0);
  \node[sub] at (2.0,3.2) {kleines Quadtree-Quadrat};
  % Die wahre Nullstelle
  \node[root] at (2.3,1.9) {};
  \node[red!80!black] at (2.65,1.95) {$\alpha_*$};
  % Newton-Iterationspunkte (Konvergenz auf alpha_*)
  \node[newtonpt,label={[above=-1pt]:$z_0$}] (z0) at (1.4,2.7) {};
  \node[newtonpt,label={[above=-1pt]:$z_1$}] (z1) at (1.8,2.35) {};
  \node[newtonpt,label={[above=-1pt]:$z_2$}] (z2) at (2.15,2.05) {};
  \node[newtonpt] (z3) at (2.28,1.92) {};
  % Pfeile
  \draw[->,teal!60!black,thick] (z0) -- (z1);
  \draw[->,teal!60!black,thick] (z1) -- (z2);
  \draw[->,teal!60!black,thick] (z2) -- (z3);
  %
  \node[sub] at (2.0,0.6) {quadratische Konvergenz};
  \node[sub] at (2.0,0.25) {gegen $\alpha_*$};
\end{scope}
%

\end{tikzpicture}
\end{document}

